\section{相关工作}
\label{sec:related-zh}

\subsection{计算机辅助作曲与后调性表示}

计算机辅助作曲的历史早于数据驱动音乐生成。OpenMusic 在可视化编程环境中表示音乐对象及其变换，使作曲家能够形式化并修订作曲过程，而不必把任务简化为音频合成 \citep{assayag1999openmusic}。music21 工具包则为符号乐谱分析、变换和交换提供程序化支持 \citep{cuthbert2010music21}。这些系统表明，符号结构的可检查性是作曲软件中的实际需求。

约束编程为显式音乐控制提供了另一项先例。\citet{anders2011constraint} 综述了以声明式约束表达和声、节奏、曲式与配器理论的系统。这一传统把音乐规则视为可检查的作曲对象，而不是语料中的隐含属性。本文解码器延续这种显式诊断取向，但使用神经模型采样的候选，而不是求解完整的约束满足问题。

音级集合理论描述移调与倒影等价关系、正常序、原型及音程类向量，用于组织非调性音高材料 \citep{forte1973structure,straus2016posttonal}。后调性理论同时区分无序集合类别与有序十二音序列、$P/R/I/RI$ 变换及十二音聚合。本文在条件表示和分析报告中保留了这一差异。

\subsection{神经符号音乐生成}

神经音乐生成涵盖不同目标、表示和控制形式 \citep{briot2020deep}。Transformer 提供标准的因果序列模型 \citep{vaswani2017attention}，Music Transformer 将自注意力用于较长的符号演奏序列 \citep{huang2019musictransformer}，后续事件表示进一步显式编码拍号与节拍结构，以生成富有表情的流行钢琴演奏 \citep{huang2020popmusic}。这些研究表明事件序列可以支持较长音乐输出，但其主要语料和评估任务并不直接测量音级集合遵循、序列变换或后调性姿态控制。

其他符号模型更关注复调协调或潜在结构。MuseGAN 使用钢琴卷帘表示生成多个器乐声部，并在摇滚伴奏场景中评估 \citep{dong2018musegan}。MusicVAE 采用分层循环解码器，在连续潜在空间中建模较长音乐序列 \citep{roberts2018hierarchical}。这些目标与多声部组织和变奏有关，但都没有提供本文任务所需的音级集合与序列形式条件。

因此，本文把 Transformer 作为已有的序列生成器使用，而不宣称架构创新。由于输出目标是量化、可编辑的乐谱片段，事件词表有意小于面向演奏的编码。音高、序列顺序、节奏、姿态、乐器和跨度条件作为前缀写入同一符号序列。

\subsection{约束感知生成}

已有多种系统支持神经生成的可控性。DeepBach 在伪 Gibbs 采样中允许对音符、节奏和终止式施加位置约束 \citep{hadjeres2017deepbach}。Anticipation-RNN 通过后向约束网络和前向标记网络把未来的一元约束纳入序列生成 \citep{hadjeres2020anticipation}。在纯神经采样之外，MorpheuS 使用优化方法保留反复模式并跟随目标张力轮廓 \citep{herremans2019morpheus}。这些方法说明显式控制具有可行性，但所处理的音乐领域与约束类型不同。

FIGARO 是与本文更接近的可解释条件生成先例。该系统利用学习描述和专家描述重构符号音乐，其中包含随时间变化的高层特征 \citep{vonrutte2023figaro}。本文把控制词汇收窄到后调性作曲材料，并在采样后加入确定性分析。实验检验在保持训练检查点和测试请求不变时，按音级、序列、节奏、姿态和音域诊断对多个候选评分是否会改变约束遵循程度。

本文方法把概率生成与确定性乐谱分析分开。因果 Transformer 先采样多个候选，后调性函数再依据音级集合覆盖率、音程向量距离、序列顺序、聚合完成率、节奏轮廓、姿态特征和乐器音域对候选排序。多候选计算增加了推断成本，但其诊断结果可以展示给作曲者并随乐谱保存。

\subsection{本文定位}

本文连接三种已有实践：计算机辅助作曲中的形式符号操作、基于事件的神经序列生成，以及约束感知选择。其具体贡献是一个面向 4--16 小节后调性草图请求的工作流程，同时返回 MusicXML 和约束遵循解释。合法的合成语料使条件--目标关系可以复现。实验在固定检查点和首个采样候选的前提下比较候选选择，并通过分别训练的前缀字段移除模型，在未改变的目标上检验各类条件。
